% =============================================================================
\section{Conclusion}
\label{sec:conclusion}
% =============================================================================

Ce projet a consisté à concevoir une métaheuristique pour l'optimisation de
trajectoires par courbes de Bézier en présence d'obstacles, sous une contrainte
de 60~secondes et une évaluation en moyenne sur 10~exécutions.

\paragraph{Phase exploratoire.}
L'implémentation et la comparaison de 19~algorithmes et variantes ont permis
d'identifier CMA-ES/BIPOP comme la méthode la plus performante sur les
instances corrélées, et de découvrir le rôle décisif des bornes étendues pour
l'instance labyrinthique \texttt{prob4}.

\paragraph{Algorithme final.}
Ces observations ont conduit à la conception d'OptiPath, une métaheuristique
basée sur CMA-ES combinant :
\begin{itemize}[nosep]
  \item un \textbf{seeding massif} ($\sim$500~solutions) avec optimisation
        locale pour identifier les bassins d'attraction prometteurs ;
  \item un \textbf{portfolio à deux marges} (5 et 15~unités) permettant de
        traiter simultanément les instances simples et labyrinthiques ;
  \item un \textbf{restart adaptatif} offrant plusieurs tentatives
        indépendantes pour trouver le bassin optimal sur les instances
        multimodales.
\end{itemize}

\paragraph{Résultats.}
OptiPath obtient en moyenne sur 10~exécutions : $36{,}80$ sur \texttt{prob1},
$31{,}42$ sur \texttt{prob2}, $29{,}02$ sur \texttt{prob3}, et $94{,}60$ sur
\texttt{prob4}.  Le gain sur \texttt{prob4} est particulièrement significatif :
un facteur~16 par rapport aux meilleures méthodes de la phase exploratoire,
grâce à l'exploitation de la vitesse paramétrique des courbes de Bézier via les
bornes étendues à marge~15.

\paragraph{Enseignements.}
Ce projet illustre l'importance de l'analyse du problème avant le choix de
l'algorithme.  La découverte du mécanisme des bornes étendues — permettant aux
points de contrôle de se placer hors du domaine pour que la courbe traverse les
obstacles à grande vitesse paramétrique — a eu un impact bien supérieur à
l'exploration de nouvelles métaheuristiques.  De même, le mécanisme de restart
adaptatif, motivé par la contrainte d'évaluation en moyenne, a transformé un
algorithme performant mais inconsistant en une solution robuste.
